\section{总结与展望}
\label{sec:summary}

在本工作中，我们在大动量有效理论框架下于 MILC 格点组态上计算了 CS 核。与我们先前的研究~\cite{LatticeParton:2020uhz}相比，本研究采用了一圈匹配 kernel，并使用了多个强子动量来提取 CS 核。我们发现，在小 $b_{\perp}$ 区域，我们的结果与微扰 QCD 一致。
在大 $b_{\perp}$ 区域，我们的结果在不确定度范围内似乎与文献中其他格点计算相一致。

在今后的研究中，我们需要
使用具有多个晶格间距的格点组态来理解有限的
晶格间距效应。我们将
采用与海夸克一致的价夸克质量以减小非幺正性效应。这类效应之一
可能是介子
波函数的虚部——目前它与微扰计算的
结果看起来不一致。显然，所有这些探索都将
耗费更多的计算资源。

